% ===== §6 =====
\section{Power as the Ontology of Culture}
\label{sec:power}

\noindent
The fourth tradition makes the strongest claim in the field, and it is placed last among the settled traditions because the antithesis will find it the hardest to answer. Its thesis is not the mild one that culture is influenced by power, which no one denies, but the radical one that power is the very being of culture: that culture, in its constitution and not merely in its uses, \emph{is} the encoding of a power order, so that to analyze a culture just is to analyze the relations of domination it congeals and transmits. On this view there is no culture outside power, no cultural form that is not, at bottom, a sedimentation of force, and the appearance of a culture that stands apart from power---an innocent web of meaning, a free play of taste, a private world between two people---is exactly the illusion by which power secures itself.

Gramsci gave the thesis its foundational form with the concept of \emph{hegemony}. The rule of a dominant class, he saw, does not rest on coercion alone; it rests, more deeply and more durably, on the manufacture of \emph{consent}, achieved in the domain of culture, where the worldview of the ruling class is diffused through the institutions of civil society until it is received by the subordinate classes as simple common sense, the natural order of things rather than the interested construction it is \citep{gramsci1971}. Culture, in this account, is the very medium in which domination is made to feel like nature, and its political function is not incidental to it but constitutive: a culture is a settlement of hegemony, a way of organizing consent to an order of power. Foucault radicalized the thesis by dissolving its remaining reference to a class that holds power. Power, for him, is not possessed but exercised, not concentrated in a ruler or a class but capillary, distributed through the finest grain of social life, and inseparable from knowledge: every ordering of knowledge is an exercise of power, every regime of power produces its knowledge, and the two are so bound that he wrote them with a single hyphenated name, \emph{power/knowledge} \citep{foucault1980pk}. Discourse, on this account, does not describe a world so much as constitute the objects it appears to describe, and the subject itself, far from being the author of culture, is one of culture's products, formed by the disciplinary and discursive practices that traverse it \citep{foucault1977dp}. Stuart Hall carried the thesis into the analysis of contemporary culture, holding culture to be not a thing to appreciate but a critical site of social action where power relations are established and may be unsettled, and showing, in his account of encoding and decoding, how the meanings a dominant order encodes into its messages may be received, negotiated, or opposed by those who decode them \citep{hall1980encoding}. Hall's is the most dialectical version, for it reopens, within the domain of cultural power, a space of contestation the others tend to close; but it remains within the tradition's founding thesis, that the field on which the contest occurs is a field of power through and through.

\subsection*{The evidence of the tradition, read from the Chinese material}

The thesis is not an abstraction, and it can be seen with particular clarity in the cultural forms this series knows best. The Confucian order of \emph{jun--chen--fu--zi}, ruler and minister, father and son, is a culture in the fullest sense, a whole texture of rite, sentiment, and meaning; and it is, at the same time and by the same texture, the encoding of a hierarchy of power, a distribution of authority naturalized into the order of the cosmos and the affections. The culture of divinity and the culture of imperial authority perform the same operation at the largest scale: the power of heaven and the power of the throne are written into the culture as the way things simply are, so that \emph{tianming}, the mandate of heaven, installs a political order in the position of the ultimate authority, the place this series has learned from psychoanalysis to call the big Other. Read through this tradition, a vast amount of what any culture contains discloses itself as the sedimentation of a power order into the seemingly natural forms of meaning and feeling, and this reading is not false. It is one of the most powerful instruments of cultural analysis there is, and the antithesis of this paper will not pretend otherwise.

\subsection*{The weld: the phallic signifier}

Before the tradition's radical consequence is drawn, one connection must be made, because it joins this tradition to the psychoanalytic line and shows that the two are not parallel tracks but a single one. If culture is a symbolic phenomenon, as the interpretive tradition holds, and if power is internal to culture, as this tradition holds, then the question arises how power is lodged \emph{within} the symbolic, and psychoanalysis gives the answer in the concept of the master signifier, and at the deepest level the \emph{phallic signifier}. In the Lacanian account the symbolic order is not a flat field of equivalent signs; it is organized around privileged signifiers that anchor its differences and orient its axis of authority, and the phallic signifier is the name for that which organizes the symbolic order's structure of power, the point around which the difference that founds hierarchy is arranged \citep{lacan2006ecrits}. This is the weld between the two lines this paper carries. Power does not enter culture from outside, as an external force bending a neutral medium; it is internal to the symbolic order's very constitution, lodged there through the phallic organization of the signifying chain, which is the same apparatus by which, as the interpretive section began to show, drive is coded and the subject's lack is opened. The line of power and the line of drive are therefore one line, joined at the phallic signifier: the coding of drive and the encoding of power are the same operation of the symbolic order seen from two sides.

\subsection*{The radical consequence, and the threat to this paper}

From the thesis, so welded, follows the consequence that makes this tradition the hardest the antithesis must face. If power is internal to the constitution of culture as such, and if there is no cultural form that is not a sedimentation of power, then there is \emph{no outside}. There is no site from which a culture free of power could be generated, no position of pure critique standing clear of the relations of domination, because the very medium in which any culture or any critique would have to be made is already, constitutively, a medium of power. A certain strain of psychoanalytically informed critique states the consequence at its sharpest: there is no pure feminism, and more generally no pure emancipatory culture, not because emancipation is not worth wanting but because there exists no system of production and social relation that is not already traversed by power, and hence no cultural form, however oppositional, that is generated in a space power has not already structured. This is the reason the tradition is placed at the head of the antithesis's difficulties. This paper means to claim that intimate relation can generate a culture of its own, endogenously, and this tradition replies, before the claim is even made, that there is no endogenous space to generate it in, that the very intimacy in which the couple would build their private world is already shot through with the phallic organization of power, that the private is a fold of the public and not an exception to it. The claim of an endogenous intimate culture is, on this tradition's terms, a naivety, and the synthesis of this paper cannot proceed until the objection has been met head-on. It will not be met by denial. It will be met, when the time comes, by conceding that power is indeed internal to all culture and turning that very concession into the ground of a different answer; but that answer belongs to the synthesis, and the tradition is left here, at the close of the thesis, in possession of its full and undiminished force.
